\documentclass[aspectratio=169,10pt,spanish]{beamer}

\input{../../../_preambulo_beamer.tex}
\input{../../../_comandos_beamer.tex}

\selectlanguage{spanish}

\title{MA1001B - Modelación Estadística para la Toma de Decisiones}
\subtitle{Lección 33: Fundamentos de pruebas de hipótesis}
\author{}
\institute{Tecnol\'ogico de Monterrey}
\date{}

\begin{document}

\begin{frame}[plain]
  \vspace{1.2cm}
  {\Large\bfseries MA1001B - Modelación Estadística para la Toma de Decisiones\par}
  \vspace{0.7cm}
  {\large Lección 33: Pruebas de hipótesis: $H_{0}$, $H_{1}$, errores de tipo I y tipo II\par}
  \vspace{0.7cm}
  {\color{TecRojo}\rule{\textwidth}{1.2pt}}
  \vspace{0.8cm}

  Facultad MA1001B\\[0.3cm]
  Periodo\\[0.3cm]
  Tecnol\'ogico de Monterrey
\end{frame}

\begin{frame}{La pregunta del lenguaje de pruebas}
  \begin{alertblock}{Pregunta guía}
    ¿Cómo se convierten $H_{0}$, $H_{1}$ y $\alpha$ en una regla de rechazo, y
    qué significan los errores de tipo I y tipo II?
  \end{alertblock}

  \vspace{0.6em}
  Al final de esta lección, usted puede:
  \begin{itemize}
    \item enunciar $H_{0}:\mu=50$ frente a $H_{1}:\mu>50$;
    \item convertir $\alpha=0.05$ en un corte de $\bar{x}$;
    \item interpretar un error de tipo I como rechazar una $H_{0}$ verdadera;
    \item explicar por qué no rechazar $H_{0}$ no prueba que $H_{0}$ sea verdadera.
  \end{itemize}

  \vfill
  \scriptsize Todos los tiempos de ciclo usados hoy son sintéticos. No se usan datos reales de empresa.
\end{frame}

\begin{frame}{Una prueba unilateral de media}
  \small
  \begin{center}
    \begin{tabular}{lr}
      \toprule
      \textbf{Cantidad} & \textbf{Valor} \\
      \midrule
      $H_{0}$ & $\mu=50$ minutos \\
      $H_{1}$ & $\mu>50$ minutos \\
      $n$, $\sigma$ conocida & 36, 8 minutos \\
      $\alpha$ & 0.05 \\
      \bottomrule
    \end{tabular}
  \end{center}

  \vspace{0.5em}
  \begin{exampleblock}{Corte}
    \[
      \mathrm{EE}=\frac{8}{\sqrt{36}}=1.333333,\qquad
      z^{*}=1.644854,
    \]
    \[
      \bar{x}_{\text{crit}}=50+1.644854\times 1.333333=52.193138.
    \]
  \end{exampleblock}
\end{frame}

\begin{frame}[fragile]{Paso 1: Hipótesis y alfa}
  \begin{columns}[T]
    \begin{column}{0.42\textwidth}
      \small
      \begin{lstlisting}
z_crit = stats.norm.ppf(0.95)
xbar_crit = (
    mu0 + z_crit * se
)
      \end{lstlisting}

      \begin{block}{Salida verificada}
        $z^{*}=1.644854$\\
        Corte $=52.193138$
      \end{block}
    \end{column}
    \begin{column}{0.58\textwidth}
      \centering
      \includegraphics[width=\textwidth]{../figuras/fundamentos_ph_01_hipotesis_alfa.png}
      \vspace{0.2em}
      \scriptsize Región de rechazo del Paso 1
    \end{column}
  \end{columns}
\end{frame}

\begin{frame}{Error de tipo I: rechazar una $H_{0}$ verdadera}
  \small
  Un error de tipo I ocurre cuando $H_{0}$ es verdadera y la muestra aún cae
  en la región de rechazo. La tasa de largo plazo de ese evento es $\alpha$.

  \vspace{0.5em}
  Simulación con semilla 42 de $10{,}000$ muestras de $\mu=50$:
  \[
    531 \text{ rechazos},\qquad
    \widehat{\alpha}=0.053100.
  \]
  La tasa simulada se sitúa cerca del $0.05$ nominal. No es idéntica a
  $0.05$ en un Monte Carlo finito.
\end{frame}

\begin{frame}[fragile]{Paso 2: Tasa simulada de tipo I}
  \begin{columns}[T]
    \begin{column}{0.40\textwidth}
      \small
      \begin{lstlisting}
xbars = samples.mean(axis=1)
rate = np.mean(
    xbars >= xbar_crit
)
      \end{lstlisting}

      \begin{block}{Salida verificada}
        Rechazos $=531$\\
        Tasa de tipo I $=0.053100$
      \end{block}
    \end{column}
    \begin{column}{0.60\textwidth}
      \centering
      \includegraphics[width=\textwidth]{../figuras/fundamentos_ph_02_tasa_tipo_i.png}
      \vspace{0.2em}
      \scriptsize Distribución muestral bajo $H_{0}$ del Paso 2
    \end{column}
  \end{columns}
\end{frame}

\begin{frame}{Paso 3: No rechazar no es una prueba}
  \begin{columns}[T]
    \begin{column}{0.42\textwidth}
      \small
      Potencia en $\mu=53$: $0.724100$.

      \vspace{0.4em}
      Tasa de tipo II: $0.275900$.

      \vspace{0.4em}
      Muestra cercana $\mu=51$:
      $\bar{x}=51.590798<52.193138$.
      No se rechaza $H_{0}$.

      \vspace{0.5em}
      \textbf{Límite:} $H_{1}$ puede ser verdadera y la prueba aún puede no
      rechazar. Ese no rechazo no prueba que $\mu=50$.
    \end{column}
    \begin{column}{0.58\textwidth}
      \centering
      \includegraphics[width=\textwidth]{../figuras/fundamentos_ph_03_tipo_ii_no_prueba.png}
      \vspace{0.2em}
      \scriptsize Potencia y tipo II del Paso 3
    \end{column}
  \end{columns}
\end{frame}

\begin{frame}{Verificación de conceptos}
  \begin{enumerate}
    \item Escriba $H_{0}$ y $H_{1}$ para esta prueba de tiempo de empaque.
    \item ¿Qué es un error de tipo I en palabras?
    \item ¿Por qué es aceptable aquí una tasa simulada de tipo I de $0.053100$?
    \item Una muestra cercana $\mu=51$ da $\bar{x}=51.59$. ¿Por qué eso no
      prueba que $\mu=50$?
  \end{enumerate}

  \vspace{0.6em}
  \begin{block}{Conexión con la Lección 34}
    La sesión puede continuar directamente con la prueba z de una muestra: el
    estadístico $z$, el valor $p$ y el supuesto de sigma conocida.
  \end{block}

  \vfill
  \scriptsize Fuente: OpenStax, \textit{Introductory Business Statistics 2e},
  Capítulo 9, Secciones 9.1--9.2. CC BY-NC-SA.
\end{frame}

\appendix



\end{document}
